% 本文件由 LaTeX 排版 Agent 根据有效 translation.json 自动生成。
% author=Gregory Thaumaturgus; work_id=0602; title=训道篇释义（圣额我略·显灵迹者）

\chapter{\WorkTitle{训道篇释义（圣额我略·显灵迹者）}}

% structure: p0001 source=h1 target=omitted reason=page-title-already-rendered-by-parent

% structure: p0002 source=h2 target=section reason=relative-heading-rank; base=h2

\section{第一章}

说这话的是\Person{所罗门}，\Person{达味}王和先知之子，他向天主的全体教会发言，他是最受尊崇的王子，是超越万人、最有智慧的先知。人的事业多么空虚和徒劳，一切占据人的追求尽皆如此！因为没有一个人能说出依附于那些事物的任何益处，这些事物是地上爬行的人竭力用身体和灵魂去追求的，同时始终服侍于短暂者，不愿用灵魂高贵的眼睛思念任何天上之事。人的生命日渐消逝，日复一日，在小时和年岁的周期中，在太阳固定的行程里，一些人不断到来，另一些人不断离去。这事如同急流的奔涌，它们带着巨响落入浩瀚无垠的深海。天主为人类所设立的一切事物都保持不变：例如，人从土而生，又归于土；大地本身稳固不移；太阳完美地绕其运行，又转回同一标记；风同样如此，奔流入海的大河以及拍打海面的微风，全都行动而不迫使大海超出其界限，它们自身也不违反所定的律法。这些事确实为了我们今生的益处而适宜地如此确立。但人所发明的事物，无论是言语还是行为，都没有限度。言语繁多，但随意和愚昧的谈论毫无益处。人类的本性在说话和听人说话两方面都贪得无厌；而且人的习惯是用闲散的眼睛观看所发生的一切。今后能发生什么，或人能做出什么尚未完成的事？有什么值得一提的新事未曾有过经验？因为我认为没有什么是可以称为新的，或思考之下会发现在古人看来是奇异或未知的。正如从前的事被遗忘而埋没，现在存在的事也将在时间中完全从后来者的知识中消失。我说这些话并非草率，仿佛现在扮演说教者。但这一切在我受托治理\Place{耶路撒冷}的希伯来人王国时，曾仔细思量过。我勤勉考察，审慎思考地上万物的本性，发现它极其纷繁；\Emph{我看见}，人注定在地上劳苦，被各种不同的辛劳所驱使，而他的工作毫无结果。此下一切事物充满怪异和可憎的精神，以致如今无人能挽回；甚至无人能想象何等彻底的虚妄已占据了所有人类事务。因为曾有一次，我与自己交谈，以为那时我在这一点上比所有在我之前的人更聪明，我精通理解比喻和万物的本性。但我认识到，我把自己投入到这样的追求毫无益处，并且如果智慧伴随着知识，那么烦恼也伴随着智慧。\par

% structure: p0004 source=h2 target=section reason=relative-heading-rank; base=h2

\section{第二章}

因此，我判断此事确是如此，便决定转向另一位生活的领航者，将自己交付于享乐，体验各种欢愉。而今我知晓这一切皆为虚空；我遏制了轻率奔放的欢笑，依中道之规约束了享乐，并对其满怀烈怒。当我察觉灵魂能制止身体耽于醉饮，且节制使情欲臣服时，便热切观察：在今生，人面前究竟摆着什么真正有价值、真正卓越的目标，使人得以达致？因为我遍历了那些被视为最尊贵的事物，诸如建造高楼、栽种葡萄园，此外还有开辟游乐庭园、获取并栽培各类果树；其中也建造了储水的大池，分布以充分灌溉树木。我也为自己蓄养了许多仆役，有男有女；有些从外邦购得，有些则是我家中所生，为我所有并役使。成群的四足牲畜，无论是牛是羊，数量超过古时任何人所拥有的，都成了我的产业。金银财宝涌入我处；我使列国的君王成为我的附庸与进贡者。我又训练了众多男女歌咏队，以和谐悦耳的歌唱愉悦我。我设宴欢饮；为此享乐，我获选了无数男女侍酒者——在这些事上，我远超耶路撒冷先前治理我的人。于是智慧在我这里衰微，而邪欲的权势却增长。因我将自己交付给眼目的一切诱惑，以及从四面八方袭来的心之暴烈情欲，并投降于享乐所给的希望，我也使我的意志成为一切可怜欢愉的奴仆。我的判断竟沦落至此，以为这些事为善，且以为我应投身其中。最终，我醒悟过来，重获视力，方知我手中所行的事全然有罪、极其邪恶，是出于不善之灵的行为。如今，世人选择的一切对象，在我看来没有一样值得赞许，或为正义之心所渴慕。因此，我既思忖智慧的益处与愚昧的祸患，便有理由大大钦佩那人：他原在轻率中奔行，而后止步，回归正道与本分。因智慧与愚昧相距甚远，彼此不同，如同白昼与黑夜。所以，选择美德者，犹如能明察一切、仰望高处、在最清晰的光明中持守其道路的人；而陷入邪恶者，犹如在无月的黑夜中无助漂泊的盲人，因黑暗而被剥夺了对事物的视力。当我思量这两种生活方式的结局时，发现后者毫无益处；我若与愚昧人为伴，则必收取愚昧的报偿。那些思想有何益处？多言有何利益？愚昧的言谈犹如从愚昧之泉涌出的溪流。再者，智慧人与愚昧人之间毫无共通之处，无论就人的纪念或是天主的报偿而言。至于人的一切事务，当它们似乎刚刚开始之时，结局便骤然临到。然而智慧人绝不与愚昧人同归于一。那时我也憎恶自己一生消耗于虚空的日子，以及我沉浸于尘世焦虑的心所度过的时光。简言之，我的一切事务皆是以劳苦和痛苦完成的，是轻率冲动的努力；另一个人，或许是智慧人或愚昧人，将继承它们——即我劳苦的冰冷果实。但当我与这些事割断联系、弃绝它们时，那人面前真正的善便向我显现——即智慧的认知与男子气概之德的拥有。若有人忽视这些，反为其他事物所炽热，这样的人便是以恶替善，追逐恶事而非美善，追逐烦恼而非平安；因为他被各种困扰所分散，日夜承载不断的焦虑，身体的重劳与心灵的无尽挂虑——他的心因那些占据他的怪异而荒谬的事务而持续动荡。因为完美的善并不在于吃喝，尽管人的滋养确实来自天主；凡为我们维持生命所赐之物，没有一样能脱离祂的眷顾而存在。但善人从天主获得智慧，也获得天上的享乐；而恶人受天主降下的灾祸击打，为情欲之病所折磨，劳苦积聚许多，并急于在天主面前使那蒙天主尊荣的人蒙羞，献上无用的礼物，将自己可怜灵魂的追逐置于虚妄与欺骗之上。\par

% structure: p0006 source=h2 target=section reason=relative-heading-rank; base=h2

\section{第三章}

因为现今的时期充满了彼此极端对立的事物——出生与死亡、植物的生长与拔除、医治与杀戮、房屋的建造与拆毁、哭泣与欢笑、哀悼与舞蹈。此刻一个人收集地上的出产，另一刻又将其抛弃；有时他热切渴望女人的\Emph{美}，有时又憎恶它。现在他寻找某种东西，然后又失去它；现在他保存，现在又抛弃；有时他杀人，有时又被杀；他说话，又沉默；他爱，又恨。因为人的事务一时处于战争状态，一时又处于和平状态；他们的命运如此无常，以至于从看似美好迅速转变为公认的灾祸。因此，让我们停止徒劳的劳苦。因为在我看来，这一切事物都像是用它们毒害的刺使人疯狂。而\OriginalTerm{不虔敬的观测时令者}{ungodly observer of the times and seasons}对这个世界张大了口，过度努力去摧毁\OriginalTerm{天主的肖像}{image of God}，如同一个从起初到末了选择与之抗争的人。因此我深信，人最大的益处乃是喜乐和行善，而这短暂的、我们唯一可能享有的快乐，只来自天主，只要正义指引我们的行为。至于天主所坚定确立的那些永恒不朽的事物，既不能从中取走什么，也不能添加什么。对一般人来说，那些事物确实可畏而奇妙；那些已经存在的事物就如此存留；那些将要存在的事物，就祂的先见而言，已经存在。此外，受损害的人有天主作他的帮助。我看见下界有\OriginalTerm{惩罚的深坑}{pit of punishment}收纳不虔敬的人，但另一个地方为虔诚的人预留。我心里想，在天主那里，万物都被判定为平等；义人和不义的人，有理性与无理性的对象，在祂的审判中都是一样的。因为他们的时间被平均分配给所有人，死亡悬在他们头上，\Emph{在这事上}，兽类和人的族类在天主的审判中相同，彼此只在于有清晰的言语上不同；其他一切事情发生在他们身上都一样，死亡平等地接纳一切，在其他种类的受造物中不比在人中更多。因为他们都有同样\Emph{生命}的气息，人并不更多；但一切，一言以蔽之，都是虚妄，从同一泥土中取得他们目前的状态，注定要消亡，并再次归于同一泥土。因为关于人的灵魂，它们是否向上飞翔，是不确定的；关于其他无理性的受造物所拥有的灵魂，它们是否向下坠落，也是不确定的。在我看来，除了快乐和享受当下的事物，没有别的益处。因为我认为一个人一旦尝过死亡，就不可能再回来享受这些事物。\par

% structure: p0008 source=h2 target=section reason=relative-heading-rank; base=h2

\section{第四章}

我抛开所有这些思虑，转而憎恶人间所行的一切压迫；有些人因此受伤哭泣哀叹，因无人保护他们而遭暴力击倒，那些本应无论如何安慰他们苦难的人也不见踪影。那些以强权为公义的人被高举到高处，然而他们也要从那里坠落。是的，在不义而胆大妄为的人中，已死的人比活着的人更好。比这两者更好的是那尚未出生的人，因为他尚未触及人间盛行的邪恶。我也看清了从邻人而来的嫉妒何等巨大，如同恶灵的刺；而\Emph{我看见了}那接受嫉妒并把它纳入胸怀的人，除了啃噬自己的心、撕裂它、消耗灵魂和身体之外，别无他事，在别人的好运中寻得无法安慰的烦恼。智慧的人宁愿选择一只手满握，若能安逸宁静，胜过双手劳作并受叛逆之灵的恶行。此外，还有一件事我知道因人的恶意而发生，与理不合。那人完全孤独，既无兄弟也无儿子，却拥有大量财产，活在贪得无厌的精神中，拒绝以任何方式向善。因此，我很想问这样的人，他为何如此劳碌，仓促逃避一切善行，却被各种求利的激情所困扰？远胜于这些的是那些采纳共同生活秩序的人，他们从中收获最好的祝福。因为当两个人以正确的精神献身于同一目标时，即使其中一人遭遇不幸，他仍因同伴在身边而得到不小的安慰。人之最大的不幸是在恶运中缺乏朋友来帮助和鼓励他。那些共同生活的人，既倍增降临的好运，也减轻逆境风暴的压力；因此白天他们以彼此坦诚的信任著称，夜里则表现为愉悦。而独居者过着一种对自己充满恐惧的生活；他没有意识到，如果一个人攻击紧密联结的人，他是在采取鲁莽危险的行动，而且三股绳不易折断。此外，我认为一个穷苦的年轻人，若是明智的，胜过没有智慧的年老王子，后者从未想过人可能从监狱被提升到王位，而那曾不义行使权力的人日后必被正义地驱逐。因为那些受制于一个有见识的青年的人——我指的是他的长辈——可能会免于烦恼。此外，后来出生的人无法赞美他们未曾经验过的另一个人，他们被不理智的判断和相反精神的冲动所引导。但在履行训道者的职务时，你要牢记这一点：使自己的生命正直，并为愚昧的人祈祷，使他们获得理解，知道如何避开恶人的行为。\par

% structure: p0010 source=h2 target=section reason=relative-heading-rank; base=h2

\section{第五章}

况且，少用舌头、说话时保持平静而正直的心，是件好事。因为将愚昧荒谬的事或心里所起的一切都用言语表达出来，是不对的；我们应当知道并反思：虽然我们远离天庭，我们说话时却在天主耳中，而且说话不冒犯我们是有益的。因为正如多种挂虑产生许多异梦和异象，愚昧的言语也与愚昧相连。再者，你要注意：凡许愿所承诺的，要确实实现。愚昧人还有一点，就是不可靠。但你当言而有信，知道与其许愿或承诺做某事而不能履行，远不如不许不诺。你务要避免污言秽语，因为天主必听见。凡以此为务的人，所得益处不过是看到自己的作为被天主挫败。因为梦多虚妄，言语也多虚妄。然而，敬畏天主是人的救恩，虽然这敬畏少见。因此，你见穷人受欺压，法官曲解律法，不必惊奇。但你要避免显出胜过掌权者。即便果真如此，因临到你身上的可怕祸患，邪恶本身也不能救你。诚然，如暴力所得之财是最有害且邪恶的产业，贪财之人永不能满足他的贪欲，也不能从邻人得善意，即便他积聚了极多财富。因为这也是虚妄。但良善使持守它的人大大喜乐，使他们强壮，赐予他们洞察万事的能力。不为这些挂虑所缠也是大事：穷人，即便为奴隶，不能饱腹，也至少享受睡眠的甜美；而贪财却带来不眠之夜和内心的焦虑。那么，还有什么比用许多焦虑和劳苦积聚财富，并嫉妒地守护它，而这一切不过是为自己维持无数祸患的根源，更愚昧的事呢？况且，这些财富终必朽坏失落，无论得者有无子女；\BibleRef{约 20:20}人自己也不情愿地注定要死，归回尘土，正与他当初受造的情形相同。\BibleRef{约 1:21}他这样注定空手离世的事实，会使这恶更令他痛苦，因为他没有考虑生命的终点与起点相似，他的劳苦无益，宁可说是为风劳力，而非为自身真实利益进展，在极不圣洁的私欲和理性之外的情欲中，并伴随着烦恼和痛苦，虚度一生。简言之，这样人的日子是黑暗，生活是悲哀。然而，这本身是好的，决不可轻视。因为人能怀着喜乐的心收获劳动的果实，拥有天主所赐而非强夺的产业，这是天主的恩赐。这样的人既不为烦恼所困，也多半不成为恶念的奴隶；他藉善行度量生命，凡事充满信心，因天主的恩赐而喜乐。\par

% structure: p0012 source=h2 target=section reason=relative-heading-rank; base=h2

\section{第六章}

此外，我要用言语指出人类中最常流行的不幸。虽然天主可以赐给人一切合他心意的事物，不剥夺任何可能迎合他欲望的对象，无论是财富、尊荣或其他使人分心之物；然而，人虽万事亨通，仿佛从天而降的唯一恶事就是无福享受，却可能仅为他人积财，对自己或邻人毫无益处而死去。我认为这是极恶的明证和清楚标记。那人无愧地被称为众多儿女之父，活了长久岁月，却在这漫长时期内心未充满善，又未曾经历死亡——这样的人，我不羡慕他的多子多孙或长寿；反而我要说，从妇人胎中堕落的流产比他还好。因为那流产既在虚妄中到来，也在遗忘中悄然离去，未尝生命之苦，未见太阳。这恶比恶人不认识善还要轻，即使他活了几千年。两者的结局都是死亡。愚昧人最显明的标志是任何欲望都不能使他满足。但谨慎的人不被这些情欲所俘。然而，多半而言，正直的生活引向贫困。好奇眼目所见扰乱许多人，点燃他们的心，以虚浮的炫耀之欲引导他们追求虚空。此外，现存之事已为人所知，且显明人不能与高于他的人争竞。确实，虚妄之事在人间流转，只增加那些以此为务之人的愚昧。\par

% structure: p0014 source=h2 target=section reason=relative-heading-rank; base=h2

\section{第七章}

因为即使一个人通过知道此生中按他心意将要发生的一切而并无太大受益（我们假定有这样的事），然而，由于人类多管闲事的活动，他仍设法窥探并表面了解人死后将要发生的事。此外，美名令人心愉悦，\BibleRef{箴 22:1} 更胜于油润身体；生命的终结优于出生，哀悼比狂欢更可取，与悲伤者同在胜于与醉酒者同在。因为事实如此：临终之人不再关心周围的光亮。审慎的愤怒胜于欢笑；因为面容的严肃使灵魂保持正直。智者的灵魂确实忧伤沮丧，但愚人的灵魂却得意洋洋，放纵于欢乐。然而，接受一个智者的责备，远比成为一整个合唱团的无用可怜之人歌声的听众更为可取。因为愚人的笑声如同许多荆棘在烈火中燃烧的爆裂声。这也是痛苦，是的，最大的邪恶，即压迫；因为它阴谋反对智者的灵魂，并试图摧毁善人所追求的崇高生活方式。此外，应当称赞的不是开始发言的人，而是结束发言的人；适度应当为自己赢得心灵的认可，而非膨胀和虚夸。再者，人必须克制愤怒，不让自己被鲁莽卷入怒气，愤怒的奴隶是愚人。此外，那些断言我们前人被赐予了更好的生活方式的人是错误的，他们未能看到智慧远非仅仅是财富的丰富，智慧比财富更为璀璨，正如银子比其影子更明亮。因为人的生命之卓越不在于获取易朽的财富，而在于智慧。谁能，请告诉我，宣告天主的眷顾，如此伟大而仁慈？谁能回忆起那些似乎被天主放过的事物？在我虚浮的往日，我考量了一切事物，\Emph{看见}一个义人持守他的正义，至死不离，却因此受害，而一个恶人与其邪恶一同灭亡。此外，义人不应显得过分正义，也不应极其过度地聪明，以免他好像在犯某种过失时，\Emph{似乎}多次犯罪。不要胆大鲁莽，以免不合时宜的死亡突然降临。抓住天主，并藉着住在他内而一无所犯，是万善中的至善。因为用不洁的手触摸纯洁之物是可憎的。但那敬畏天主而顺服自己的人，逃避一切相反之事。智慧在帮助上的效力胜过城中一队最强有力的人，并且它也常常公义地赦免那些失职的人。因为没有人不跌倒。你也绝不可听从不敬虔之人的言语，以免你亲耳听到针对自己的话，如邪恶仆人愚昧的谈话，因而心中刺痛，随后自己在许多行为中转而咒骂。这一切我都知道，因为我从天主接受了智慧，之后却失去了，再也不能恢复原样。因为智慧离我而去，到了无限遥远之处，进入无底的深渊，以致我再也无法抓住它。因此之后我完全放弃寻求它；我不再考虑不敬虔之人的愚昧和虚妄计谋，以及他们疲惫、烦乱的生活。怀着这样的心态，我被推向这些事物本身；被一种致命的激情抓住，我认识了女人——她像网罗或类似的东西。因为她的心诱骗那些经过她的人；她只要手接手，就牢牢抓住一个人，仿佛用链子拖着他。要从她手中解脱，唯有找到天主作为仁慈而警觉的监督；因为被罪奴役的人无法（否则）逃脱它的掌握。此外，在所有女人中，我寻找她们应有的贞洁，却一个也没找到。确实，人在一千人中可能找到一个贞洁的男人，但女人从未有。最重要的是，我观察到：人受天主创造时心思单纯，但他们为自己招致种种推理和无穷的疑问，并自称寻求智慧，却在虚浮的言语中浪费生命。\par

% structure: p0016 source=h2 target=section reason=relative-heading-rank; base=h2

\section{第八章}

再者，智慧若存于一人身上，也显现在拥有者的面容上，使他的面容发光；反之，厚颜无耻一旦占据人心，人一被看见，就被判定为可憎。因此，完全有必要留意君王的话语，并以各种方式避免发誓，尤其是以天主之名发誓。同时，注意恶言或许是适宜的，但必须防备任何对天主的亵渎。因为当祂施加惩罚时，不可能指责祂，也不可能反驳唯一之主、唯一之君的谕令。但人若遵守神圣诫命，远离恶人的言语，则更好、更有益。因为智慧人预知并辨别适时来到的审判，看出它将是正义的。因为人类生活中的一切都等候从上而来的报应；但恶人似乎真不知道，既然有强大的眷顾照管着他，将来的事没有一样可隐藏。他确实不知道将要发生的事，因为没有人能适当地向他宣告其中之一；没有谁可强大到能阻止夺取他生命的天使，也没有什么方法可以消除那定期的死亡。但正如在战争中被俘的人，只看见四面被截断的逃跑之路，同样，人的一切不虔敬也一同全然消灭。我每次思想人们图谋伤害邻舍的诸事是何等之多、何等之大，就不禁惊异。但我深知，不虔敬的人因自身沉迷于虚幻，而被从今生夺去，并被除掉。因为天主的眷顾审判因祂的极大含忍，并不迅速临到众人，恶人犯罪后也并未立即受罚——因此他以为可以更加犯罪，似乎可以免于惩罚，不理解即使经过很长时间，违法者也逃不过天主的知觉。此外，至善就是敬畏天主；因为不虔敬的人一旦离开祂，就不会被容许长久滥用其愚昧。但有一种极其邪恶虚假的意见常常在人间流行，关于义人和不义之人。因为人们对二者形成违反真相的判断；真正正义的人得不到应有的认可，而不虔敬的人却被视为谨慎正直。我认为这是最严重的错误之一。曾几何时，我真的以为至善在于吃喝，凡在今生能极尽享乐的人就是最得天主宠眷的；我幻想这种享乐是人生唯一的慰藉。于是，我只致力于这个念头，昼夜不辍于一切曾被发现能给予人奢华欢愉的事。我由此学到的是：沉浸于这些事的人，无论如何劳苦，也绝不能找到真正的好。\par

% structure: p0018 source=h2 target=section reason=relative-heading-rank; base=h2

\section{第九章}

那时我曾以为，所有人都被判决同样的事。若有什么智者实行正义，远离不义，并因审慎避免众人的仇恨（这确实是悦乐天主的事），在我看来，此人却是徒劳。因为义人与不虔敬者，善人与恶人，洁净与不洁，敬拜天主与不敬拜的，似乎都有同一个结局。因为不义的人与善人、发假誓的人与完全避免起誓的人，我都怀疑他们趋向同一结局，一种邪恶的见解偷偷潜入我心，即众人都以相似的方式终结。但现在我知道这是愚人的想法，是错误与欺骗。他们大肆宣称：死了的人完全灭亡，活人比死人更可取，即使他躺在黑暗中，像狗一样度过人生旅程，\Emph{这}至少比死狮还好。因为活人至少知道他们将要死；但死人什么都不知道，在他们完成必走的旅程之后，也没有什么酬报为他们预备。而且死人的仇恨与爱情都终止了；因为他们的嫉妒已经消失，他们的生命也熄灭了。凡是一去不返的人，在什么事上都没有分。错误仍弹着这样的琴弦，也提供这样的忠告：人哪，你是什么意思，为何不娇养自己，用各种美味填饱自己，用酒灌满自己？你岂不看见这些事是天主赐给我们，让我们无节制地享受的吗？穿上刚洗过的衣服，用没药膏抹你的头，看看这个女人那个女人，虚度你虚妄的一生。因为除了这个，无论是今世或死后，都没有什么留给你。但要利用一切机会；因为没有人会为这些事向你算账，人在人所行的事在人类圈子之外也完全无人知晓。至于阴府\OriginalTerm{阴府}{Hades}，无论它是什么，据说我们要往那里去，它既没有智慧，也没有明智。这些是虚无之人所说的话。但我确切知道：那些似乎跑得最快的人不能完成那伟大的赛跑；那些在人的判断中被视为有能可畏的人，也不能在那可怕的争战中得胜。同样，智慧也不由面包丰裕证明，明智也不惯于与财富为伍。我也不祝贺那些以为众人都会遭遇同样事情的人。但确实，那些沉浸于这些想法的人在我看来是睡着了，没有考虑到他们像鱼和鸟一样突然被捉住，将被祸患吞灭，很快遇到他们应得的报应。我也如此珍视智慧，以至于我宁愿把一座小而居民少的城市，即使被一位大君王率军围攻，也视为伟大而有能力，只要城中有一个智慧的人，无论多么贫穷。因为这样的人能够拯救他的城市脱离敌人和壁垒。而其他人，也许不认识那个贫穷的智慧人；但对我来说，我大大偏爱智慧中蕴含的力量，胜过单纯民众的势力。然而在这里，智慧与贫穷共居，被轻视。但将来，它要说话，声音比寻求邪恶的王子与暴君更有权威。因为智慧比铁更强；而一个人的愚昧给许多人带来危险，即使他被许多人鄙视。\par

% structure: p0020 source=h2 target=section reason=relative-heading-rank; base=h2

\section{第十章}

再者，苍蝇落入没药中，窒息在其中，使那馨香的膏油和涂抹这膏油的动作都成了可憎的事；同时记念智慧与愚昧，实属不当。智慧人确实是自己善行的引领者；但愚昧人倾向错误的道路，决不会以他的愚昧作为追求高尚的向导。是的，他的思念也是虚空，充满愚昧。但是，我的朋友，如果有敌对的精神临到你，要勇敢抵挡它，因为知道天主能平息甚至大堆的过犯。这些也是那一切邪恶的元首和父亲的行为：愚昧人被高举，而富有智慧的人被贬抑；罪的奴仆骑马奔驰，而献身于天主的人不光彩地步行，恶人却洋洋得意。但如果有人谋害别人，他忘了自己首先且独自为自己预备了网罗。破坏他人安全的人，必被蛇咬伤。搬石头的人，必受不小的劳苦；劈柴的人，必在自己的武器上携带危险。如果斧头从柄上脱落，从事这工作的人必陷入麻烦，徒劳地聚集，并不得不使用更多他那邪恶而短促的力量。再者，蛇的咬伤是隐蔽的；念咒的人不能平息疼痛，因为他们是虚妄的。但善人既为自己也为邻人做善事；而愚昧人必因他的愚昧沉入灭亡。一旦他开口，就开始愚昧地说话，很快便结束，处处显露出他的无知。此外，人不可能知道任何事，也不能从人那里知道从起初是什么，或将来是什么。因为谁能宣告这些呢？此外，那不知道去往好城的人，在眼目和整个面容上承受邪恶。我预言那城中君王是少年、首领是贪食者的城有祸了。但我称那君王是自由人之子的好地为有福：在那里，被托付治理权的人将在适当的时候收获善果。但懒惰和闲散的人成为嘲笑者，使房屋衰败；为满足自己的贪食欲而滥用一切，如同金钱的奴隶，为了微小的代价，他们甘心做一切卑贱和可鄙的事。服从君王、统治者或掌权者也是应当的，不要对他们怀恨，也不要说任何冒犯的话。因为暗中说的话总有风险可能被公开。因为快而带翅膀的使者将一切传达给那唯一富足大能的君王，在其中履行一种既属神又合理的服务。\par

% structure: p0022 source=h2 target=section reason=relative-heading-rank; base=h2

\section{第十一章}

再者，将你的饼和那些维持人生命所需之物施舍给（穷人）是正义的事。因为虽然你似乎即刻将其浪费在某些人身上，如同将饼撒在水面上，但随着时间的推移，你的仁慈将被视为对你有益。也要慷慨地施舍，将你的财物分给许多人；因为你不知道明日将带来什么。再者，云彩不会留住丰沛的雨水，而是将阵雨倾泻在地上。树也不能永远站立；即使人饶过它，也终会被风吹倒。但许多人也渴望预先知道天上将要发生的事；曾有人审视云彩、等候风，却不去收割和簸扬，信赖虚妄，完全不能知道任何可能从天主而来的未来之事；正如人不能告诉怀胎的妇人生下什么一样。要在季节播种，待收割的时候来到，就收获你的果实。因为不能显明在所有自然事物中什么比那些更好。巴不得一切事情都顺利！实在，当人思量太阳是好的，此生是甜蜜的，有多年可不断享受是愉快的，而死亡是恐怖、无尽的邪恶、使我们归于无有的事物，他就认为应当享受当下一切明显的快乐。他也给年轻人这样的劝告：要充分利用青春的时期，放纵心意于各类享乐，满足情欲，行自己眼中看为好的事，注视那使人喜悦的，避开那不喜悦的。但我要对这样的人说：\Quote{愚拙的人啊，我的朋友}，因为你没有期待天主将要临到这一切的审判。纵欲和放荡是邪恶的，我们身体污秽的淫乱带来死亡。因为愚昧伴随青春，而愚昧引向灭亡。\par

% structure: p0024 source=h2 target=section reason=relative-heading-rank; base=h2

\section{第十二章}

此外，你应当趁年轻时就敬畏天主，在将自己交给恶事之前，在主的伟大而可畏的日子来临之前，那时太阳不再发光，月亮和众星也不再发光，而是在那万物的风暴和动荡中，上面的权势——即守护世界的天使——将要被震动；以致勇士们将衰败，妇女们将停止劳作，逃入他们住所的黑暗之处，并关闭所有的门。妇女因恐惧而不再推磨，发出最微弱的声音，如同最小的鸟；所有不洁的妇女将沉入地下；城市及其染血的政府将等待从上而来的报复，而最苦涩和最血腥的时期像盛开的杏花悬在他们头上，连续的惩罚像成群的飞蝗即将来临，违反者像黑色可鄙的刺山柑植物被抛出路旁。善人将欢欢喜喜地离去，往自己的永久居所；但卑劣者将用哀哭充满他们的一切地方，积存的银子或炼净的金子也不再有用处。因为一场重击将落在万物之上，甚至及于井旁的水罐，和可能留在洼处的器皿轮子，当时期的进程终结，如水的生命中那带着洗涤的时期过去之时。对于躺在地上的人，只有一个救恩，就是他们的灵魂认识并飞向那创造他们的天主。那么，我再说一遍我已经说过的话：人的境况全是虚妄，没有什么能超越附着于人的发明之物上的极端虚妄。我谨慎地宣讲的劳动是多余的，因为我是在试图教导一个民族，他们如此不愿接受教导或医治。实在，高尚的人才是理解智慧之言所需要的。再者，我虽然已经年迈，度过了漫长的一生，仍努力通过真理的奥秘找出那些中悦天主的事。我知道，智者的训诫使心灵活跃和受激发，不亚于身体在受到刺棒或钉子钉入时惯常的反应。有些人将从一位良善的牧者和导师那里领受这些智慧课程，然后似乎众口同声、同心合意地更详尽地阐述交托给他们的真理。但多言无益。我也不劝你，我的朋友，写下关于合宜之事的虚妄之言，从中除了厌倦的劳苦一无所获。最后，我需要用这样一个结论：你们哪，人，看，我现在明确而简短地嘱咐你们：敬畏天主，祂是主，也是万有的监察者，并遵守祂的诫命；并且相信将来一切都要分别受审判，各人要照自己的行为，或善或恶，得到公义的报应。\par

\SourceRecord{Translated by S.D.F. Salmond.；Ante-Nicene Fathers；Vol. 6.；Edited by Alexander Roberts, James Donaldson, and A. Cleveland Coxe.；Buffalo, NY: Christian Literature Publishing Co.,；1886.；<http://www.newadvent.org/fathers/0602.htm>.}
